\chapter{Praktische Umsetzung: Implementierung}
\label{sec:implementierung}

\section{Einleitung}
In diesem Kapitel wird die technische Umsetzung des \textit{multimodalen Sprach-Assistenten} beschrieben. Der Fokus liegt auf der Integration der einzelnen Komponenten der KI-Pipeline, bestehend aus Speech-to-Text, Natural Language Processing sowie der visuellen Ausgabe der analysierten Inhalte.

Ziel der Implementierung ist die Realisierung eines durchgängigen Systems, das gesprochene Sprache verarbeitet, strukturiert analysiert und in visuelle Informationen überführt.

\section{Workflow der Implementierung}
Die Entwicklung erfolgte modular entlang der einzelnen Stufen der KI-Pipeline:

\begin{enumerate}
    \item \textbf{Audioeingabe:} Aufnahme oder Upload von Sprachdaten über die Benutzeroberfläche (Streamlit).
    
    \item \textbf{Speech-to-Text Verarbeitung:} Umwandlung der Audiodaten in Text mittels \textbf{OpenAI Whisper}.
    
    \item \textbf{NLP-Verarbeitung:} Analyse des transkribierten Textes mit \textbf{spaCy}, inklusive Tokenisierung und Part-of-Speech Tagging.
    
    \item \textbf{Informationsextraktion:} Identifikation relevanter sprachlicher Elemente, insbesondere Verben und Nomen.
    
    \item \textbf{Visuelle Ausgabe:} Nutzung extrahierter Nomen als Suchbegriffe für eine Bilddatenbank (z.\,B. Unsplash API).
    
    \item \textbf{UI-Darstellung:} Visualisierung der Ergebnisse innerhalb einer Streamlit-Webanwendung.
\end{enumerate}

\section{Integration der Systemkomponenten}

\subsection{Speech-to-Text Modul}
Das Speech-to-Text Modul basiert auf dem Whisper-Modell, welches Audiosignale in Textsequenzen umwandelt. Die Implementierung erfolgt so, dass sowohl Audio-Dateien als auch potenziell Live-Eingaben verarbeitet werden können.

\subsection{NLP-Analyse Pipeline}
Die Verarbeitung des Textes erfolgt über spaCy. Dabei wird der Text zunächst tokenisiert und anschließend grammatikalisch analysiert.

Die wichtigsten Ergebnisse sind:
\begin{itemize}
    \item Identifikation von Verben (Handlungen)
    \item Extraktion von Nomen (Objekte)
    \item Strukturierung des Satzes in linguistische Einheiten
\end{itemize}

\subsection{Informationsextraktion und Mapping}
Die extrahierten Nomen werden als semantische Schlüsselbegriffe genutzt. Diese werden in Suchanfragen transformiert, um passende visuelle Inhalte zu finden.

Dieser Schritt stellt eine einfache Form der semantischen Abbildung zwischen Sprache und Bild dar.

\subsection{Benutzeroberfläche (Streamlit)}
Die gesamte Interaktion erfolgt über eine Streamlit-Webanwendung. Diese übernimmt:

\begin{itemize}
    \item Texteingabe und Audio-Upload
    \item Darstellung der NLP-Ergebnisse
    \item Visualisierung der markierten Verben im Text
    \item Anzeige der generierten bzw. abgerufenen Bilder
\end{itemize}

\section{Praktische Umsetzungsschritte}
Die Implementierung wurde in einzelne Python-Module unterteilt:

\begin{itemize}
    \item \texttt{speech\_to\_text.py}: Verarbeitung von Audioeingaben mit Whisper
    \item \texttt{nlp\_analysis.py}: Textanalyse mit spaCy (Tokenisierung, POS-Tagging)
    \item \texttt{image\_generator.py}: Generierung von Bild-URLs basierend auf Nomen
    \item \texttt{app.py}: Streamlit-Frontend und Orchestrierung der Pipeline
\end{itemize}

\section{Besondere Herausforderungen und Lösungen}

\begin{itemize}
    \item \textbf{Modellabhängigkeiten:} Die Integration von Whisper erforderte eine stabile Python-Umgebung mit kompatiblen Versionen von NumPy und PyTorch.
    
    \item \textbf{Audioverarbeitung:} Unterschiedliche Audioformate wurden durch temporäre Konvertierung in ein einheitliches Format gelöst.
    
    \item \textbf{Semantische Genauigkeit:} Die Qualität der Bildzuordnung hängt stark von der Auswahl relevanter Nomen ab, weshalb eine Filterung der Token implementiert wurde.
\end{itemize}

\section{Testing und Qualitätssicherung}
Die Qualitätssicherung erfolgte durch manuelle und funktionale Tests der Pipeline:

\begin{itemize}
    \item Test verschiedener Sprachinputs (einfach, komplex, mehrsatzig)
    \item Validierung der korrekten POS-Erkennung durch spaCy
    \item Überprüfung der Bildzuordnung anhand semantischer Konsistenz
    \item UI-Tests innerhalb der Streamlit-Anwendung
\end{itemize}

\section{Zusammenfassung}
Die Implementierung zeigt die erfolgreiche Realisierung einer modularen KI-Pipeline. Durch die klare Trennung der einzelnen Verarbeitungsschritte konnte ein flexibles und erweiterbares System geschaffen werden, das Spracheingaben in strukturierte und visuelle Informationen überführt.